%%=============================================================================
%% CAPITULO 4 -- APRESENTACAO, ANALISE E DISCUSSAO DOS DADOS
%% Versao dummy: resumo dos capitulos 05-resultados e 06-discussao da
%% dissertacao real.
%%=============================================================================

\section{Análise e discussão dos dados}

O módulo de cálculo, escrito em Python puro, foi conferido contra o gabarito
do conjunto de casos de cálculo, com o critério de aceite, fixado antes da
conferência, de uma diferença de no máximo R\$~0,01 em cada um dos 600 casos.
O módulo cumpriu o critério em todos eles, e os 17 testes automatizados,
que incluem os casos-limite de prazo zero, de 29 e 30 dias na fronteira do
IOF e de taxa de rendimento nula, passam sem acesso à rede. A base regulatória
foi dividida em 3.989 trechos, e, nas perguntas de teste, a busca híbrida teve
os maiores valores em todas as métricas, exceto no recall@1, em que a busca
semântica ficou à frente.

Esses resultados são parciais, porque a avaliação técnica automatizada das
respostas do agente ainda depende da execução do experimento. Entre as
limitações, a clareza das respostas é validada por LLM-as-a-judge e por
métrica de legibilidade, sem validação humana complementar, e 10 das 51
perguntas de teste da base ficaram sem trecho correto entre os 10 primeiros
resultados. A discussão também reposiciona o sentido de democratização: para o
perfil Sem Reservas, sem capacidade real de investimento, a resposta correta
do sistema é não recomendar produto algum, em respeito à vedação da Resolução
CVM nº~30/2021 de recomendar produto inadequado ao perfil do cliente.
